\chapter{Experiments and Results}
\label{chap:experiments}

This chapter presents the experimental evaluation of multitask learning approaches for forest attribute estimation from Sentinel-1 SAR imagery. Two datasets are used to investigate different aspects of multitask learning: TreeSatAI-CHM provides a well-recognized benchmark for exploring fundamental MTL challenges like learning rate mismatches and gradient conflicts, while iMAESTRO offers comprehensive reference data suitable for testing physics-aware deep learning through allometric constraints. The experiments address how shared architectures learn task-specific and common features, where gradient conflicts emerge, and whether multitask learning improves performance over single-task baselines. Models are implemented in PyTorch and trained on NVIDIA A100 GPUs.

\section{TreeSatAI-CHM Dataset Experiments}
\label{sec:treesat_experiments}

The TreeSatAI-CHM dataset combines the established TreeSatAI benchmark with generated canopy height maps. While the small patch size ($6\times 6$ pixels) and limited reference data constrain absolute performance, this dataset serves as a controlled environment for understanding multitask learning dynamics. The experiments focus on qualitative and quantitative analysis of how classification and regression tasks interact in a shared encoder. All experiments use the dataset described in Section~\ref{sec:treesat_chm_dataset} with consistent train/validation/test splits.

\subsection{Single-Task Baseline Experiments}
\label{subsec:treesat_baselines}

\subsubsection{Experimental Setup}

This section presents the experimental evaluation of single-task baseline models for tree species classification and canopy height estimation. These baselines establish reference performance levels against which multitask learning approaches are compared.

\begin{itemize}
    \item Batch size: 128 (large batch for stable gradient given small $6\times 6$ patch size)
    \item Optimizer: AdamW with weight decay $10^{-5}$ for regularization
    \item Early stopping: patience of 5 epochs monitoring validation loss
    \item Random seed: 42 for reproducibility
    \item Epochs: 100
\end{itemize}

\paragraph{Classification Loss} The classification task uses BCEWithLogitsLoss (binary cross-entropy with logits), treating the problem as multi-label classification where each patch can belong to multiple genus classes. This approach aligns with the TreeSatAI benchmark methodology. The loss is weighted by class frequencies to address substantial class imbalance in the dataset. The weighting scheme follows inverse square root frequency normalization, balancing between uniform weighting (which ignores class imbalance) and inverse frequency weighting (which can overemphasize rare classes).

\paragraph{Regression Loss} The regression task uses root mean squared error (RMSE) as the loss function. A masking mechanism excludes pixels with invalid CHM values (below \SI{1}{\meter} threshold) from loss computation, as these represent non-forest areas or data gaps.

\begin{itemize}
    \item \textbf{MLP models:} Hidden dimensions (512, 512, 512), dropout 0.2
    \begin{itemize}
        \item MLP Classifier: learning rate $10^{-5}$
        \item MLP Regressor: learning rate $5\times 10^{-5}$
    \end{itemize}
    \item \textbf{U-Net models:} Base channels $B = 256$, dropout 0.2
    \begin{itemize}
        \item U-Net Classifier: learning rate $10^{-5}$
        \item U-Net Regressor: learning rate $5\times 10^{-5}$
    \end{itemize}
\end{itemize}

For U-Net models, the large channel capacity ($B = 256$, resulting in 1024 channels at bottleneck) was chosen to compensate for limited spatial extent of small patches, allowing the network to encode sufficient information in the channel dimension when spatial dimensions are heavily compressed.

\subsubsection{Classification Results}

Table~\ref{tab:baseline_classification} presents the classification performance of MLP and U-Net baselines on the test set.

\begin{table}[htbp]
\centering
\caption{Baseline Classification Performance}
\label{tab:baseline_classification}
\footnotesize
\begin{tabular}{l|c|cccc|cccc}
\hline
 & & \multicolumn{4}{c|}{\textbf{Weighted}} & \multicolumn{4}{c}{\textbf{Micro}} \\
\cline{3-10}
\textbf{Model} & \textbf{OA} & \textbf{Prec.} & \textbf{Rec.} & \textbf{$F_1$} & \textbf{mAP} & \textbf{Prec.} & \textbf{Rec.} & \textbf{$F_1$} & \textbf{mAP} \\
\hline
TreeSatAI Bench & -- & 33.29 & 7.13 & 10.09 & 29.42 & 63.01 & 7.13 & 12.82 & 33.09 \\
MLP Classifier & 7.10 & 33.55 & 8.26 & 10.61 & 29.13 & 58.26 & 8.26 & 14.47 & 32.65 \\
U-Net Classifier & 9.38 & 37.12 & 12.82 & 17.14 & 26.71 & 45.14 & 12.82 & 19.96 & 29.02 \\
\hline
\end{tabular}
\end{table}

The MLP classifier achieves 7.10\% overall accuracy on the 15-class genus classification task. This result aligns with TreeSatAI benchmark performance for MLP models on Sentinel-1 data, which reported similarly modest accuracy for SAR-based classification. The high precision (58.26\%) combined with low recall (8.26\%) indicates the model is conservative in its predictions, correctly identifying a subset of samples but missing the majority of positive instances.

The U-Net classifier shows improvement over the MLP baseline, achieving 9.38\% overall accuracy and 19.96\% micro F1 score. The improvement in recall (from 8.26\% to 12.82\%) suggests that the convolutional architecture captures spatial patterns that help identify more positive instances. However, this comes with reduced precision (from 58.26\% to 45.14\%), indicating a shift toward more aggressive predictions.

The relatively low mAP scores (around 27--33\%) reflect the difficulty of the task:, that distinguishing tree genera from SAR backscatter alone is challenging because many genera share similar structural properties and thus similar radar signatures.

\subsubsection{Regression Results}

Table~\ref{tab:baseline_regression} presents the canopy height regression performance.

\begin{table}[htbp]
\centering
\caption{Baseline Regression Performance}
\label{tab:baseline_regression}
\begin{tabular}{l|cc|cc|cc}
\hline
 & \multicolumn{2}{c|}{\textbf{Train}} & \multicolumn{2}{c|}{\textbf{Val}} & \multicolumn{2}{c}{\textbf{Test}} \\
\cline{2-7}
\textbf{Model} & \textbf{RMSE (m)} & \textbf{$R^2$} & \textbf{RMSE (m)} & \textbf{$R^2$} & \textbf{RMSE} & \textbf{$R^2$} \\
\hline
MLP Regressor & 7.23 & 0.052 & 7.27 & 0.020 & 7.27 & 0.030 \\
U-Net Regressor & 7.24 & 0.051 & 7.20 & 0.037 & 7.22 & 0.042 \\
\hline
\end{tabular}
\end{table}

Both models achieve similar RMSE values around 7.2--7.3~m, with the U-Net showing slight improvement on validation and test sets. The $R^2$ values are low (around 0.02--0.05), indicating that the models explain only a small fraction of the variance in canopy height.

The U-Net regressor achieves marginally better test performance (RMSE 7.22~m, $R^2$ 0.042) compared to the MLP regressor (RMSE 7.27~m, $R^2$ 0.030). This is still a modest improvement, but could suggest that preserving spatial structure through the encoder-decoder architecture with skip connections provides useful information for height estimation that is lost when the input is flattened.

\subsubsection{Analysis}

\paragraph{1) MLP Performance Aligns with TreeSatAI Benchmark} The MLP classification results are consistent with the original TreeSatAI benchmark, which reported that Sentinel-1 SAR data alone provides limited discriminative power for tree species classification compared to optical imagery. The benchmark noted that MLP models on Sentinel-1 achieve substantially lower accuracy than ResNet models on Sentinel-2 or aerial imagery, primarily because SAR backscatter patterns lack the spectral diversity present in multispectral optical data.

\paragraph{2) U-Net Preserves Spatial Information} The U-Net models outperform MLP baselines on both tasks, though the improvements are modest. For classification, the U-Net achieves approximately 2 percentage points higher accuracy and 5 percentage points higher F1 score. For regression, the improvement is smaller but consistent across validation and test sets. These results suggest that, even though the $6\times 6$ patches are relatively small, the spatial patterns within them still contain useful information that the convolutional architecture can exploit, whereas the MLP loses this structure by flattening the input.

\paragraph{3) Data Quality Validation} To verify that the limited performance is due to the inherent difficulty of SAR-based prediction rather than data quality issues, additional experiments were conducted using Sentinel-2 optical imagery on the same dataset. The Sentinel-2 experiments achieved 32.8\% classification accuracy and 5.78~m RMSE for height regression. These substantially better results confirm that the dataset contains learnable patterns relating imagery to forest parameters, and that the modest SAR performance reflects the information content of radar backscatter rather than problems with the data pipeline or labels.

\paragraph{4) Implications for Multitask Learning} The U-Net architecture provides a reasonable foundation for multitask learning, as it already demonstrates the ability to extract useful features for both classification and regression. The shared encoder in the multitask architecture may benefit from the regularization effect of optimizing for multiple objectives, potentially improving generalization compared to single-task training.

\subsection{Multitask Learning Experiments}
\label{subsec:treesat_mtl}

\subsubsection{Experimental Setup}

\paragraph{Loss Scale Mismatch and Training Challenges}

Initial experiments revealed substantial challenges in balancing the classification and regression objectives. The two losses operate on fundamentally different scales: the weighted binary cross-entropy loss for classification typically ranges from 0.3 to 0.8, while the RMSE for regression spans 5 to 15~m. This magnitude difference creates severe optimization difficulties.

Preliminary experiments explored fixed weight combinations ranging from $w_{\text{cls}}:w_{\text{reg}}$ ratios of 1:1 to 0.3:200. Even with per-task loss normalization, training exhibited evident imbalances. A learning rate suitable for one task often proved inappropriate for the other: for instance, a learning rate of $5\times 10^{-4}$ that enabled stable regression training caused the classification head to diverge or overfit rapidly, while rates small enough for stable classification training (e.g., $10^{-5}$) were insufficient for meaningful regression learning.

\begin{itemize}
    \item Base channels: $B = 256$ (matching U-Net regressor baseline)
    \item Dropout: 0.2 in regression decoder
    \item Classification head: same class weighting scheme (inverse square root frequency, normalized) as single-task classifier
    \item Batch size: 128
    \item Epochs: 100 (early stopping)
\end{itemize}

To address these challenges, the experiments used task-specific learning rates combined with cosine annealing scheduling. The optimizer uses separate parameter groups for the shared encoder, classification head, and regression decoder, each with independent learning rates:

\begin{itemize}
    \item Shared encoder: $5\times 10^{-6}$
    \item Classification head: $5\times 10^{-6}$
    \item Regression decoder: $5\times 10^{-5}$
    \item Uncertainty parameters: 0.0001 (log learning rate)
\end{itemize}

A cosine annealing scheduler gradually reduces learning rates over training, with minimum learning rate set to $10^{-6}$. Uncertainty weighting is enabled, allowing the model to automatically balance task contributions.

\subsubsection{Results and Analysis}

Table~\ref{tab:mtl_comparison} presents the performance comparison between single-task baselines and the multitask model.

\begin{table}[htbp]
\centering
\caption{Multitask Learning Performance Comparison}
\label{tab:mtl_comparison}
\begin{tabular}{lccccccc}
\hline
\textbf{Model} & \textbf{OA} & \textbf{Prec.} & \textbf{$F_1$} & \textbf{Rec.} & \textbf{mAP} & \textbf{RMSE} & \textbf{$R^2$} \\
\hline
U-Net Classifier & 9.38\% & 45.14\% & 19.96\% & 12.82\% & 29.02\% & -- & -- \\
U-Net Regressor & -- & -- & -- & -- & -- & 7.27 & 0.042 \\
MultiTaskUNet & 7.02\% & 59.20\% & 14.24\% & 8.09\% & 32.99\% & 7.26 & 0.045 \\
\hline
\end{tabular}
\end{table}

The multitask model achieves 7.02\% overall accuracy with 14.24\% F1 and 32.99\% mAP for classification, while maintaining 7.26~m RMSE and 0.045 $R^2$ for regression. Compared to single-task baselines, the classification overall accuracy decreases slightly (from 9.38\% to 7.02\%). The regression performance remains comparable to the baseline with slightly improved $R^2$ (from 0.042 to 0.045).

The results show an asymmetric relationship between the two tasks. The regression task maintains stable performance from the shared encoder, suggesting that features learned for classification also help with height estimation. This makes sense because species-specific patterns like texture and backscatter characteristics relate to forest structure, which affects canopy height. However, the classification task shows some performance drop compared to its single-task baseline.

Several factors explain this asymmetry. The shared encoder has limited capacity, and when training for both tasks at once, it must split this capacity between learning species-specific features and height-related features. Since height varies continuously within and across species, the regression task may dominate what the encoder learns, leaving less room for the detailed species distinctions needed for classification. Additionally, the classification task tends to overfit during multitask training. With only 15 genus classes and unbalanced class sizes, the classification head can quickly memorize training patterns. The regression task, which predicts continuous values at each pixel, provides stronger training signals that can overwhelm the classification objective. During training, validation classification metrics often peaked early then declined, while regression metrics kept improving.

The tasks also differ in complexity. Classification makes one prediction per $6\times 6$ patch, while regression makes 36 predictions (one per pixel). This means regression contributes much more gradient signal per sample, which can dominate encoder updates. Furthermore, species classification likely needs sensitivity to subtle backscatter differences that separate similar genera (like distinguishing \textit{Quercus} from \textit{Fagus}, both broadleaf deciduous trees), while height estimation mainly responds to overall backscatter strength and texture patterns. The shared encoder may learn features better suited for height estimation, losing the fine details needed for classification.

\begin{figure}[htbp]
    \centering
    \includegraphics[width=0.62\textwidth]{chm-samples.png}
    \hfill
    \includegraphics[width=0.32\textwidth]{chm-scatter-plot.png}
    \caption[Multitask model predictions on TreeSatAI-CHM test set]{Multitask model predictions on TreeSatAI-CHM test set. Left: Example predictions showing VV/VH input, true CHM, and predicted CHM. The model captures general height patterns but struggles with fine spatial detail. Right: Scatter plot of predicted versus true height values. The clustering around 20--30~m indicates limited sensitivity to height variations, reflecting SAR backscatter limitations.}
    \label{fig:treesat_mtl_predictions}
\end{figure}

\subsubsection{Gradient Analysis}

To understand task interactions in the shared encoder, gradient cosine similarity was computed between classification and regression gradients at three encoder depths across training epochs. Gradient cosine similarity has been widely used to analyze task conflicts in multitask learning, where negative values indicate conflicting gradients that can harm optimization \cite{yu2020gradient}. Recent work has further developed this analysis to understand when and where tasks compete for shared parameters \cite{liu2022autolambda,guo2023dynamic}.

\begin{figure}[htbp]
    \centering
    \includegraphics[width=0.7\textwidth]{gradient-encoder.png}
    \caption[Gradient cosine similarity between classification and regression tasks]{Gradient cosine similarity between classification and regression tasks at encoder layers. Positive values indicate aligned gradients where tasks agree on parameter updates; negative values indicate conflicts where tasks push parameters in opposite directions.}
    \label{fig:treesat_gradient}
\end{figure}

\paragraph{1) Early Encoder Layers (enc1)} Gradient alignment fluctuates substantially, occasionally showing negative cosine similarity. This suggests low-level feature extraction experiences fundamental conflict. Classification may benefit from edge-preserving features for distinguishing genus boundaries, while regression may benefit from smoothing features that suppress speckle noise.

\paragraph{2) Middle Encoder Layers (enc2)} More consistent positive alignment appears, indicating mid-level spatial patterns benefit both tasks. Both classification and regression require understanding of forest structure, leading to compatible feature requirements at this level.

\paragraph{3) Bottleneck} The strongest positive alignment occurs at the bottleneck, suggesting abstract semantic representations are shared. Both species identity and canopy height relate to forest structure at a conceptual level, and the bottleneck captures this shared semantic information.

The predominantly positive gradient cosine similarities indicate that classification and regression tasks generally share compatible learning objectives. The shared encoder learns features useful for both tasks, though early layers experience occasional conflicts. The gradient conflict primarily occurs in early layers while alignment strengthens in deeper layers, which is consistent with multitask learning intuition and related work. The task-specific learning rates successfully balance these competing objectives, preventing either task from dominating encoder updates.

The primary goal of the subsequent multitask experiments is to investigate whether joint training on classification and height regression can improve feature learning from SAR imagery. Even if overall performance remains modest due to the inherent limitations of SAR data, understanding the interactions between these tasks in a shared encoder, provides insight into the relationship between species identity and canopy structure as encoded in radar backscatter.

\subsection{Ablation Study: Alternative Parameter Sharing Strategies}
\label{subsec:treesat_ablation}

To better understand the baseline multitask architecture, this section compares three parameter sharing strategies: the hard-sharing MultiTaskUNet, Cross-Stitch Networks \cite{misra2016crossstitch}, and Cross-Task Attention.

\subsubsection{Qualitative Comparison of Architectures}

\begin{table}[htbp]
\centering
\caption{Architecture Comparison}
\label{tab:arch_comparison}
\begin{tabular}{lcccc}
\hline
\textbf{Architecture} & \textbf{Sharing Degree} & \textbf{Task Freedom} & \textbf{Param.} & \textbf{Complexity} \\
\hline
MultiTaskUNet & Highest & Lowest & $\sim$8M & Low \\
Cross-Stitch & Medium & High & $\sim$16M & Medium \\
Cross-Attention & High (encoder) & Medium (bottleneck) & $\sim$10M & High \\
\hline
\end{tabular}
\end{table}

\paragraph{Hard Parameter Sharing (MultiTaskUNet)}

The baseline architecture shares all encoder parameters between tasks and uses task-specific decoders. This represents the maximum degree of sharing, forcing both tasks to use identical low-level and mid-level features. The architecture is simple and parameter-efficient but provides no mechanism for tasks to learn specialized representations until the decoder branches.

\paragraph{Cross-Stitch Networks}

Cross-stitch networks \cite{misra2016crossstitch} introduce learnable linear combinations between task-specific subnetworks at multiple depths. After a shared initial encoder block, the network splits into parallel classification and regression pathways. At designated layers, cross-stitch units combine activations from both pathways:
\begin{align}
f_{\text{cls}}^{\text{new}} &= \alpha_{11} f_{\text{cls}} + \alpha_{12} f_{\text{reg}}, \\
f_{\text{reg}}^{\text{new}} &= \alpha_{21} f_{\text{cls}} + \alpha_{22} f_{\text{reg}},
\end{align}
where the $\alpha$ parameters are learnable scalars initialized to favor task-specific features (e.g., $\alpha_{11} = \alpha_{22} = 0.9$, $\alpha_{12} = \alpha_{21} = 0.1$). This architecture provides substantially more flexibility than hard sharing, allowing each task to learn specialized features while selectively incorporating useful information from the other task. However, the duplicated encoder pathways approximately double the parameter count.

The cross-stitch approach is well-suited when tasks have related but distinct optimal representations. The learned $\alpha$ values reveal which layers benefit from sharing versus specialization.

\paragraph{Cross-Task Attention}

The cross-attention variant introduces multi-head attention modules that allow each task to selectively attend to features from the other task at the bottleneck. After a shared encoder, task-specific bottleneck convolutions produce separate feature maps for classification and regression. The cross-attention module then computes:
\begin{align}
f_{\text{cls}}^{\text{attended}} &= f_{\text{cls}} + \text{Attention}(Q=f_{\text{cls}}, K=f_{\text{reg}}, V=f_{\text{reg}}), \\
f_{\text{reg}}^{\text{attended}} &= f_{\text{reg}} + \text{Attention}(Q=f_{\text{reg}}, K=f_{\text{cls}}, V=f_{\text{cls}}).
\end{align}

This design maintains a shared encoder but allows task-specific refinement of the bottleneck representation through attention. The attention mechanism can learn complex, content-dependent interactions between task features, potentially capturing non-linear relationships that cross-stitch units cannot express.

However, the attention mechanism introduces additional computational overhead and learnable parameters (query, key, value projections for each attention head). For the small patch sizes and limited training data in this study, the additional complexity may lead to overfitting rather than improved performance.

\subsubsection{Results and Analysis}

\begin{table}[htbp]
\centering
\caption{Ablation Study Results}
\label{tab:ablation_results}
\begin{tabular}{lcccc}
\hline
\textbf{Model} & \textbf{OA} & \textbf{Weighted mAP} & \textbf{Test RMSE (m)} & \textbf{Test $R^2$} \\
\hline
MultiTaskUNet & 8.15\% & 28.99\% & 7.26 & 0.045 \\
Cross-Stitch & 5.81\% & 24.52\% & 7.31 & 0.037 \\
Cross-Attention & 5.19\% & 23.18\% & 7.18 & 0.047 \\
\hline
\end{tabular}
\end{table}

The ablation results reveal distinct behavioral patterns for each architecture:

\paragraph{Cross-Stitch Networks} achieve 5.81\% overall accuracy with $R^2$ of 0.037, showing weaker classification performance compared to the baseline MultiTaskUNet but slightly better regression. The learned cross-stitch parameters likely favor the regression pathway, as the cross-stitch mechanism allows the regression branch to develop specialized features without being constrained by classification-optimal representations. This behavior is consistent with the original cross-stitch paper's finding that task-specific representations emerge when the architecture provides freedom to specialize \cite{misra2016crossstitch}. However, with the limited training data available, the nearly-doubled parameter count may contribute to overfitting, and the classification task may receive insufficient signal from the regression-optimized shared representations.

\paragraph{Cross-Task Attention} shows a similar pattern, achieving 5.19\% overall accuracy and the best regression performance ($R^2 = 0.047$). The attention mechanism at the bottleneck appears to favor regression-relevant information flow. This may occur because the regression task produces spatially-distributed features that provide richer key-value pairs for attention, while the classification task condenses all spatial information through global average pooling immediately after the bottleneck. The attention mechanism may thus learn to prioritize regression-relevant feature interactions simply because they provide more diverse attention targets. Additionally, the attention mechanism introduces substantial additional complexity that may be inappropriate for the small patch sizes and limited training data in this study. With only $6\times 6$ pixel inputs, the attention operates on minimal spatial extent, limiting its ability to learn meaningful long-range dependencies. The additional parameters may lead to overfitting on the classification task, which has fewer supervision signals per sample.

The ablation results suggest that for this particular combination of tasks, dataset size, and input dimensions, the simpler hard-sharing architecture provides a reasonable balance between tasks. The more complex architectures introduce flexibility that the limited data cannot effectively leverage, leading to task imbalance rather than mutual improvement.

These findings align with the general principle that model complexity should match data availability. Cross-stitch networks were originally demonstrated on ImageNet-scale datasets with thousands of samples per class, while this study operates with approximately 40{,}000 training samples distributed across 15 imbalanced classes. Similarly, attention mechanisms typically benefit from larger spatial extents and more training data to learn meaningful attention patterns. For future work with larger datasets or higher-resolution inputs, the cross-stitch or cross-attention architectures may prove more effective. 

\subsection{Summary}
\label{subsec:treesat_summary}

The TreeSatAI-CHM experiments establish that multitask learning can combine genus classification and canopy height regression from Sentinel-1 SAR data, though with modest absolute performance due to the limited information content of radar backscatter. The key findings are:

First, the baseline experiments confirm that both tasks are learnable from SAR data. U-Net models outperform MLP baselines on both tasks, showing that spatial patterns within even small $6\times 6$ patches contain useful information. Validation experiments with Sentinel-2 optical data (achieving 32.8\% classification accuracy and 5.78~m RMSE) confirm that the dataset quality is sound and the modest SAR performance reflects inherent radar limitations rather than data issues.

Second, the multitask model maintains regression performance while showing slight classification degradation. This asymmetry suggests that features learned for classification transfer well to height estimation, but not vice versa. The regression task's pixel-level predictions (36 per patch) provide much stronger gradient signals than classification's patch-level predictions (1 per patch), which can dominate encoder learning. Task-specific learning rates and uncertainty weighting help balance these competing objectives.

Third, gradient analysis reveals where tasks cooperate and compete. Early encoder layers show occasional gradient conflicts, suggesting low-level feature requirements differ (edge preservation for classification vs. noise smoothing for regression). Deeper layers and the bottleneck show strong positive alignment, indicating that abstract representations of forest structure benefit both tasks. This pattern informed the architectural ablation study.

Fourth, the ablation experiments with cross-stitch and cross-attention architectures show that added flexibility does not improve performance on this dataset. Both alternative architectures achieve worse classification results while maintaining similar regression performance, likely because the limited training data and small patch sizes cannot effectively leverage the additional parameters. The simpler hard-sharing architecture provides a reasonable balance for this scale of data.

\section{iMAESTRO Dataset Experiments}
\label{sec:imaestro_experiments}

The iMAESTRO dataset provides comprehensive reference data for genus segmentation, canopy height, and biomass estimation across three European forest sites. With larger patch sizes ($64\times 64$ pixels) and complete co-registration of all target variables, this dataset enables investigation of physics-aware multitask learning through allometric constraints. All experiments use the spatially-blocked train/validation/test splits described in Section~\ref{sec:imaestro_dataset}.

\subsection{Single-Task Baseline Experiments}

\subsubsection{Experimental Setup}

Three independent U-Net models are trained for genus segmentation, height regression, and biomass regression. The deeper architecture (four downsampling stages) exploits the larger patch size compared to TreeSatAI-CHM.

Training configuration:
\begin{itemize}
    \item Batch size: 16 (suitable for $64\times 64$ patches with 128 base channels)
    \item Optimizer: AdamW with weight decay $10^{-4}$
    \item Dropout: 0.2
    \item Learning rate: $10^{-4}$ for all models
    \item Epochs: 100 with early stopping (patience 10)
    \item Random seed: 42
\end{itemize}

Task-specific configurations:

\begin{itemize}
    \item \textbf{Segmentation:} Cross-entropy loss with 16 rare genera excluded (ignore index). Predicts 6 dominant genera.
    \item \textbf{Height regression:} Masked RMSE loss excluding invalid pixels (Non-forest pixels).
    \item \textbf{Biomass regression:} Masked RMSE loss with same masking strategy.
\end{itemize}

\subsubsection{Results}

Table~\ref{tab:imaestro_baselines} presents the single-task baseline performance.

\begin{table}[htbp]
\centering
\caption{iMAESTRO Single-Task Baseline Performance}
\label{tab:imaestro_baselines}
\begin{tabular}{lcccccc}
\hline
\textbf{Task} & \textbf{Metric} & \textbf{Train} & \textbf{Val} & \textbf{Test} \\
\hline
Segmentation & Overall Accuracy & 68.2\% & 64.5\% & 63.8\% \\
             & Mean IoU & 0.512 & 0.478 & 0.471 \\
Height       & RMSE (m) & 4.82 & 5.15 & 5.23 \\
             & $R^2$ & 0.312 & 0.268 & 0.255 \\
Biomass      & RMSE (t/ha) & 48.3 & 52.1 & 53.7 \\
             & $R^2$ & 0.031 & 0.026 & 0.024 \\
\hline
\end{tabular}
\end{table}

The segmentation baseline achieves 63.8\% overall accuracy on the test set with mean IoU of 0.471. This performance is substantially better than TreeSatAI-CHM classification (9.38\%), reflecting the larger patch size, simpler task (6 genera vs 15), and higher-quality reference data.

Height regression achieves test RMSE of 5.23~m with $R^2 = 0.255$. This represents meaningful predictive power, explaining approximately 25\% of height variance. The performance is better than TreeSatAI-CHM ($R^2 = 0.042$), benefiting from larger patches and higher-quality labels.

Biomass regression shows the weakest performance with test RMSE of 53.7~t/ha and $R^2 = 0.024$. The low $R^2$ reflects the label quality and the fundamental challenge of C-band SAR saturation for high biomass values. SAR backscatter saturates at approximately 100--150~t/ha for C-band, limiting sensitivity to biomass variations in mature forests.

\subsection{Multitask Learning with Allometric Constraints}

\subsubsection{Experimental Setup}

The multitask model jointly predicts genus segmentation, canopy height, and biomass using a shared encoder with three task-specific decoders. The allometric constraint couples height and biomass predictions.

Model configuration:

\begin{itemize}
    \item Base channels: 128 (same as baselines)
    \item Dropout: 0.2 for segmentation decoder, 0.1 for regression decoders
    \item Uncertainty weighting: enabled with learnable $\sigma$ parameters
    \item Allometric constraint: $\lambda_{\text{allom}} = 0.1$, $\alpha = 0.0673$, $\beta = 2.5$        
\end{itemize}

Training Configuration:

\begin{itemize}
    \item Batch size: 16
    \item Optimizer: AdamW with weight decay $10^{-4}$
    \item Learning rate: $10^{-4}$ for all parameters
    \item Epochs: 100 with early stopping (patience 10)
\end{itemize}

\subsubsection{Results and Analysis}

Table~\ref{tab:imaestro_mtl} compares multitask performance to single-task baselines.

\begin{table}[htbp]
\centering
\caption{iMAESTRO Multitask Learning Performance}
\label{tab:imaestro_mtl}
\begin{tabular}{lcccc}
\hline
\textbf{Task} & \textbf{Metric} & \textbf{Baseline} & \textbf{Multitask} & \textbf{Change} \\
\hline
Segmentation & Overall Accuracy & 63.8\% & 62.1\% & -1.7\% \\
             & Mean IoU & 0.471 & 0.458 & -2.8\% \\
Height       & RMSE (m) & 5.23 & 5.11 & -2.3\% \\
             & $R^2$ & 0.255 & 0.265 & +3.9\% \\
Biomass      & RMSE (t/ha) & 53.7 & 51.2 & -4.7\% \\
             & $R^2$ & 0.024 & 0.039 & +62.5\% \\
\hline
\end{tabular}
\end{table}

The multitask model achieves 62.1\% segmentation accuracy (vs 63.8\% baseline), showing slight degradation. Height regression improves to $R^2 = 0.265$ (+3.9\% relative), while biomass shows the largest improvement with $R^2 = 0.039$ (+62.5\% relative increase, though still low in absolute terms).

The learned uncertainty parameters converged to approximately 1.0 for segmentation, 0.4 for height, and 0.3 for biomass, reflecting relative task difficulty and loss scales.

\begin{figure}[htbp]
    \centering
    \includegraphics[width=0.9\textwidth]{mtl-all-tasks-samples.png}
    \caption[Multitask model predictions on iMAESTRO test set]{Multitask model predictions on iMAESTRO test set. Example predictions showing VV/VH input and predictions for all three tasks from the shared encoder architecture. The multitask model maintains comparable visual quality to single-task baselines while using a single shared encoder.}
    \label{fig:imaestro_mtl_samples}
\end{figure}

\begin{figure}[htbp]
    \centering
    \includegraphics[width=0.9\textwidth]{height-biomass-scatter-plot.png}
    \caption[Height-biomass relationship in multitask predictions]{Height-biomass relationship in multitask predictions. Scatter plot showing the correlation between predicted height and biomass values, demonstrating that the allometric constraint successfully maintains the ecological relationship between these variables.}
    \label{fig:imaestro_height_biomass}
\end{figure}

Height and biomass losses exhibited coupled behavior after epoch 15, with correlated fluctuations suggesting that the allometric constraint successfully linked the two regression tasks. The validation metrics for all three tasks improved steadily until around epoch 35, and after segmentation and height metrics stabilized, biomass continued to increase slightly.

\subsection{Gradient and Representation Analysis}

\subsubsection{Gradient Alignment in the Shared Encoder}

\begin{figure}[htbp]
    \centering
    \includegraphics[width=0.7\textwidth]{gradient-comparison.png}
    \caption[Gradient cosine similarity between task pairs in the shared encoder]{Gradient cosine similarity between task pairs in the shared encoder across training epochs. Positive values indicate aligned gradients where tasks agree on parameter updates; negative values indicate conflicting gradients where tasks push parameters in opposite directions.}
    \label{fig:imaestro_gradients}
\end{figure}

The gradient analysis reveals dynamic task relationships throughout training. In the first 15 epochs, all three task pairs exhibit highly variable gradient alignment with frequent sign changes, indicating tasks initially compete for shared encoder capacity. The height-biomass pair (green line) shows the strongest positive correlation during this phase, with cosine similarities frequently exceeding 0.6, suggesting these two regression tasks naturally share low-level feature requirements.

After epoch 15, gradient patterns stabilize. The segmentation-height pair (blue line) settles into moderate positive alignment (cosine similarity 0.4--0.7), indicating genus classification and height estimation benefit from similar mid-level features. The segmentation-biomass pair (orange line) shows the weakest and most variable alignment throughout training, with cosine similarities ranging from -0.3 to 0.6, suggesting genus classification and biomass estimation require partially conflicting features.

The height-biomass pair maintains strong positive alignment (0.4--0.9) after the initial training phase, particularly from epochs 20--40. This sustained gradient agreement provides direct evidence that the allometric constraint successfully couples the two regression tasks.

\subsubsection{Feature Similarity Across Decoder Layers}

\begin{figure}[htbp]
    \centering
    \includegraphics[width=\textwidth]{cka.png}
    \caption[Centered Kernel Alignment (CKA) similarity between task-specific decoder representations]{Centered Kernel Alignment (CKA) similarity between task-specific decoder representations at different layers. Each heatmap shows CKA similarity between decoder layers of two tasks. Layer 4 is closest to the bottleneck, while layer 1 is nearest to the output. Higher CKA values (red) indicate more similar representations.}
    \label{fig:imaestro_cka}
\end{figure}

CKA analysis quantifies how similarly task-specific decoders transform shared encoder features. The segmentation-height comparison (left panel) shows high similarity (CKA > 0.7) at the deepest decoder layer (layer 4), immediately after the shared bottleneck. As we move toward shallower layers (layers 3, 2, 1), CKA similarity decreases progressively, reaching 0.51--0.66 at layer 3 and 0.51--0.54 at layers 2 and 1. This gradient of decreasing similarity demonstrates that decoders gradually specialize for their respective tasks.

The segmentation-biomass comparison (middle panel) exhibits lower overall similarity, with CKA values ranging from 0.25 to 0.51. The deepest layer (layer 4) shows moderate similarity (CKA $\approx$ 0.5), substantially lower than the segmentation-height pair, confirming biomass estimation requires fundamentally different feature processing than genus classification.

The height-biomass comparison (right panel) shows the strongest similarity among all task pairs, with CKA values of 0.77--0.81 at layer 4 and remaining high (0.64--0.81) through layers 3, 2, and 1. This sustained high similarity across all decoder depths provides compelling evidence that height and biomass estimation share feature processing strategies throughout the decoder pathway.

\begin{figure}[htbp]
    \centering
    \includegraphics[width=\textwidth]{representations-across-layers.png}
    \caption[Layer-wise representation analysis across decoder depths]{Layer-wise representation analysis across decoder depths. Visualization of how task-specific representations evolve from the shared bottleneck through decoder layers, showing progressive specialization for each task while maintaining alignment for related tasks (height and biomass).}
    \label{fig:imaestro_representations}
\end{figure}

\subsection{Summary}

The iMAESTRO experiments tested multitask learning on larger patches ($64\times 64$ pixels) with three tasks: genus segmentation, canopy height regression, and biomass regression. The key experimental findings are:

First, baseline performance shows that all three tasks are learnable from SAR data, with segmentation achieving 63.8\% accuracy and height regression reaching $R^2 = 0.255$. Biomass regression shows weaker performance ($R^2 = 0.024$), reflecting both label quality issues and C-band SAR saturation at high biomass values (100-150 t/ha). The larger patch size compared to TreeSatAI-CHM provides more spatial context, leading to better absolute performance.

Second, the multitask model with allometric constraints improves both regression tasks while maintaining segmentation performance. Height $R^2$ increases from 0.255 to 0.265 (+3.7\%), and biomass $R^2$ increases from 0.024 to 0.039 (+62.5\% relative improvement). The allometric constraint (with weight $\lambda = 0.1$) couples height and biomass predictions, enforcing the ecological relationship $\log(B) = \alpha + \beta \times \log(H)$ with parameters $\alpha = 0.0673$ and $\beta = 2.5$.

Third, gradient analysis shows strong alignment between height and biomass tasks (cosine similarity 0.4-0.9 throughout training), confirming that the allometric constraint guides both tasks toward compatible feature learning. Segmentation shows weaker alignment with biomass, suggesting genus classification needs different features than continuous attribute estimation. The learned uncertainty parameters (1.0 for segmentation, 0.4 for height, 0.3 for biomass) reflect relative task difficulty and automatically balance task contributions.

Fourth, CKA analysis reveals how representations evolve through the network. The bottleneck shows high similarity across all task pairs (0.64-0.81), indicating shared abstract features. Similarity decreases through decoder layers as tasks specialize, but height-biomass similarity remains high (0.64-0.81 across all decoder layers) due to the allometric coupling. This sustained similarity provides evidence that physics-aware constraints maintain consistent representations throughout the network.
